% ============================================================
%  Scalar multiplication, double-and-add, ECDLP, one-way functions
% ============================================================

\section{Scalar multiplication}

\begin{frameddefn}[label={def:scalarmul}]{Scalar multiplication}
  For an integer $k \ge 1$ and a point $P$:
  \[
    kP = \underbrace{P + P + \cdots + P}_{k \text{ times}}, \qquad 0P = \mathcal{O}, \qquad (-k)P = k(-P).
  \]
\end{frameddefn}

The ``multiplication'' is between an \textbf{integer} and a \textbf{point}: two points are never multiplied together. The integer $k$ is called the \textit{scalar}.

\subsection{Naive method: repeated addition}

Compute $G, 2G = G + G, 3G = 2G + G, \ldots$ one after the other. This needs $k-1$ additions: perfect for printing the whole sequence on the toy curve, hopeless for real keys ($k \approx 2^{256}$ would need more additions than there are atoms in the universe).

\subsection{Fast method: double-and-add}

Write $k$ in binary, $k = (b_{m-1} \ldots b_1 b_0)_2$. Reading the bits from the most significant one, \textit{doubling} shifts the partial scalar left by one bit and \textit{adding} $P$ sets the last bit to 1.

\begin{framedalgo}[label={alg:dbladd}]{Double-and-add (left-to-right)}
  \begin{algorithmic}[1]
    \Function{ScalarMul}{$k, P$}
      \State $R \gets \mathcal{O}$
      \For{$i = m-1$ \textbf{downto} $0$}
        \State $R \gets 2R$ \Comment{double}
        \If{$b_i = 1$}
          \State $R \gets R + P$ \Comment{add}
        \EndIf
      \EndFor
      \State \Return $R$
    \EndFunction
  \end{algorithmic}
\end{framedalgo}

\begin{framedex}{$k = 6 = (110)_2$}
  \begin{center}
    \begin{tabular}{c c l l}
      \toprule
      bit & value & operation & $R$ \\
      \midrule
      $b_2$ & 1 & $2\mathcal{O} + G$ & $G$ \\
      $b_1$ & 1 & $2G + G$ & $3G$ \\
      $b_0$ & 0 & $2 \cdot 3G$ & $6G$ \\
      \bottomrule
    \end{tabular}
  \end{center}
  The first step is trivial ($R$ just becomes $G$), then we need only 2 doublings and 1 addition, instead of the 5 additions of the naive method.
\end{framedex}

\begin{framedprop}{Cost}
  Double-and-add uses at most $m$ doublings and $m$ additions, where $m = \lfloor \log_2 k \rfloor + 1$ is the number of bits of $k$. For a 256-bit key: about 256 doublings and on average 128 additions.
\end{framedprop}

There is also a right-to-left variant (process bits from $b_0$, keeping a running $2^i P$), exactly like square-and-multiply for modular exponentiation. Real implementations use \textit{constant-time} variants (e.g. Montgomery ladder) because the plain algorithm does an addition only when the bit is 1, so its running time leaks the bits of the key (\textit{timing side channel}).

\section{The Elliptic Curve Discrete Logarithm Problem}

\begin{framedprob}[label={prob:ecdlp}]{ECDLP}
  Given a curve $E(\mathbb{F}_p)$, a base point $G$ of order $n$ and a point $Q \in \abk{G}$, find the integer $k \in \set{0, \ldots, n-1}$ such that $Q = kG$.
\end{framedprob}

The name comes from the analogy with ordinary exponentiation: $kG$ is the additive version of $g^k$, and finding $k$ is a ``logarithm''.

\begin{itemize}
  \item \textbf{Brute force}: try $G, 2G, 3G, \ldots$ until you meet $Q$: $O(n)$ steps.
  \item \textbf{Baby-step giant-step} and \textbf{Pollard's rho}: $O(\sqrt{n})$ steps. These are the best known generic algorithms for well-chosen curves.
  \item For $n \approx 2^{256}$, $\sqrt n \approx 2^{128}$ operations: completely infeasible. That is what ``128-bit security'' means.
  \item \textbf{Pohlig--Hellman}: if $n$ factors into small primes, the problem splits into small problems. This is why $n$ must be a large prime.
\end{itemize}

\begin{framedobs}{Toy curve = no security}
  On our curve $n = 9$: anyone can recover $k$ from $Q$ by trying at most 8 multiples. The toy curve shows the \textit{mechanism}, not the \textit{hardness}.
\end{framedobs}

\section{One-way functions and hash functions}

\begin{frameddefn}[label={def:owf}]{One-way function}
  A function $f$ is \textbf{one-way} if it is efficiently computable, but for a random input $x$ every efficient algorithm, given $y = f(x)$, finds some $x'$ with $f(x') = y$ only with negligible probability.
\end{frameddefn}

Under the ECDLP assumption, $f(k) = kG$ is (believed to be) one-way. Actually it is more than that: it is a \textbf{bijection} between $\set{0, \ldots, n-1}$ and $\abk{G}$ and a \textbf{group homomorphism}:
\[
  f(k_1 + k_2) = (k_1 + k_2)G = k_1 G + k_2 G = f(k_1) + f(k_2).
\]
This structure is exactly what makes digital signatures (ECDSA) and key exchange (ECDH: $k_A(k_B G) = k_B(k_A G)$) possible.

\begin{frameddefn}[label={def:hash}]{Cryptographic hash function}
  A function $H : \set{0,1}^* \to \set{0,1}^\ell$ (arbitrary input length, fixed output length) with:
  \begin{itemize}
    \item \textbf{preimage resistance}: given $h$, it is hard to find $x$ with $H(x) = h$ (i.e. one-wayness);
    \item \textbf{second-preimage resistance}: given $x$, it is hard to find $x' \neq x$ with $H(x') = H(x)$;
    \item \textbf{collision resistance}: it is hard to find any pair $x \neq x'$ with $H(x) = H(x')$.
  \end{itemize}
  In practice a good hash function also behaves like a \textbf{random oracle}: outputs look uniformly random, there is no algebraic relation between $H(x)$ and $H(x')$, and flipping one input bit flips each output bit with probability $1/2$ (\textbf{avalanche effect}).
\end{frameddefn}

\begin{center}
  \small
  \begin{tabularx}{\textwidth}{X X X}
    \toprule
    Property & $k \mapsto kG$ & Hash $H(x)$ (e.g. Keccak-256) \\
    \midrule
    Deterministic, easy to compute & yes & yes \\
    Hard to invert (preimage) & yes (ECDLP) & yes \\
    Input domain & fixed: $\set{1,\ldots,n-1}$ & any length \\
    Collisions & \textbf{none} (injective on $\set{0..n-1}$) & must exist, but hard to find \\
    Algebraic structure & homomorphism: $f(a+b) = f(a)+f(b)$ & must have none \\
    Output distribution & uniform over $\abk{G}$ if $k$ uniform & should look uniform for \textit{any} inputs \\
    Avalanche effect & not required, partially present & required \\
    Used for & keys, signatures, key exchange & integrity, commitments, addresses \\
    \bottomrule
  \end{tabularx}
\end{center}
